\documentclass[11pt]{nyjm}

\usepackage{amsmath,amsthm,amssymb,amsfonts,mathtools}
\usepackage{microtype}
\usepackage{graphicx}
\usepackage{url}
\usepackage{listings}
\usepackage{xcolor}
\usepackage[colorlinks=false,hidelinks]{hyperref}

\newtheorem{theorem}{Theorem}[section]
\newtheorem{lemma}[theorem]{Lemma}
\newtheorem{proposition}[theorem]{Proposition}
\newtheorem{corollary}[theorem]{Corollary}
\newtheorem{axiom}[theorem]{Axiom}
\theoremstyle{definition}
\newtheorem{definition}[theorem]{Definition}
\newtheorem{remark}[theorem]{Remark}

\lstdefinelanguage{Lean}{
  keywords={axiom, theorem, def, structure, inductive, example, open, namespace, import, theory, Prop, Type, Nat, Real, Int, Q, Vector, Matrix, eval, HasSignature, Rle, Pos, Req, Rneg, true, by, intro, exact, let, if, then, else, in},
  sensitive=true,
  comment=[l]{--},
  morecomment=[s]{/-}{-/},
  string=[b]"
}

\lstset{
  language=Lean,
  basicstyle=\ttfamily\small,
  breaklines=true,
  keywordstyle=\color{blue},
  commentstyle=\color{green!60!black},
  stringstyle=\color{red}
}

\title{The F\textsubscript{1} Square: A Formal Research Program for the Riemann Hypothesis} 

\author[F1 Consortium]{The F1 Square Research Consortium}
\address{Crown City, OH}
\email{consortium@citizengardens.org}

\keywords{Riemann Hypothesis, F1-geometry, Hodge Index Theorem, Constructive Real Analysis, Lean 4}
\subjclass[2010]{Primary 11M26; Secondary 14G40, 11M38}

\begin{document} 
 
\begin{abstract}  
We present a formal (Lean 4) research program that reduces the Riemann Hypothesis (RH) to the construction of a single missing geometric object: the arithmetic surface $\mathbb{S} = \operatorname{Spec} \mathbb{Z} \times_{\mathbb{F}_1} \operatorname{Spec} \mathbb{Z}$ equipped with an intersection pairing satisfying a Hodge index theorem. We formalize a conditional theorem: \textit{if such a surface exists with an ample class $H$, a symmetric bilinear intersection form, and a Hodge index theorem (negative-definiteness on $H^\perp$), then all non-trivial zeros of the Riemann zeta function lie on the critical line.}

To demonstrate the framework's consistency, we instantiate the intersection theory on a finite truncation of primes (the $T_1$ set) and verify that the Gram matrix has signature $(1, 3n)$, matching the full-space Hodge index signature. Crucially, this work does \emph{not} prove RH nor construct $\mathbb{S}$. The existence of the arithmetic surface remains the sole open core. This manuscript serves as a defensive publication of the formal scaffold.
\end{abstract} 

\maketitle
\tableofcontents

\section{Introduction and Epistemic Framework}
The Riemann Hypothesis remains open. This document outlines a precise, machine-checkable reduction of RH to a problem in absolute arithmetic geometry.

\subsection{Core Philosophy}
Our formalization is built upon the \texttt{UOR-Foundation} constructive analysis substrate. It adheres to a strict honesty policy: every unproven assumption is explicitly declared as an \texttt{axiom} or marked as open. The Lean 4 codebase enforces this via a Continuous Integration audit that forbids \texttt{sorry} in the production core without explicit declaration.

\subsection{The Missing Analog}
Analogous to the function field case where the Hodge index theorem on $C \times C$ proves RH for curves over finite fields, we target the arithmetic surface $\mathbb{S}$. The known components—the characteristic-1 base, the scaling flow (Frobenius analogue), and the explicit formula as a trace—are verified. The missing piece is the 2-dimensional self-product $\mathbb{S}$ and its intrinsic intersection theory.

\section{The Sedona Spine \& Projective Hodge Index}
We formalize the intersection theory for a finite truncation of $n$ primes. The basis of divisor classes is:
\[
\{H_i, V_i, \Delta, \Gamma_j \mid i, j = 0, \dots, n-1\}
\]
where $H_i, V_i$ are the horizontal/vertical rulings, $\Delta$ is the diagonal, and $\Gamma_j$ is the scaling graph for the $j$-th prime.

\subsection{Template Structure}
In \texttt{IntersectionTemplate.lean}, we define the Gram matrix of size $3n+1$ with a block structure:

\begin{lstlisting}[caption={Template Data Structure}]
structure TemplateData (n : Nat) where
  d : Real
  a : Fin n -> Fin n -> Real
  b : Fin n -> Fin n -> Real
  c : Fin n -> Fin n -> Real
  a_sym : forall i j, Req (a i j) (a j i)
  b_sym : forall i j, Req (b i j) (b j i)
  c_sym : forall i j, Req (c i j) (c j i)
  trace_condition : Real
\end{lstlisting}

The construction exploits the sourced product-of-curves relations:
\begin{align*}
\langle H_i, H_j \rangle &= 0, \quad \langle V_i, V_j \rangle = 0, \quad \langle H_i, V_j \rangle = \delta_{ij}, \\
\langle H_i, \Delta \rangle &= 1, \quad \langle V_i, \Delta \rangle = 1, \quad \langle \Delta, \Delta \rangle = d.
\end{align*}

\subsection{Ample Class and Orthogonal Complement}
We define the ample class vector:
\[
H = \sum_{i=1}^n (H_i + V_i) + c \cdot \Delta
\]
with $c$ chosen such that $\langle H, H \rangle > 0$. The orthogonal complement $H^\perp$ is the subspace of vectors orthogonal to $H$ under the intersection form.

\section{Explicit Axioms \& Disclaimers}
The core theorems in \texttt{IntersectionTemplate.lean} rely on two fundamental linear algebra axioms. These are the only unproven assumptions in the formal conditional proof.

\subsection{Axiom 1: Orthogonal Basis Existence}
\begin{axiom}[Orthogonal Diagonalisation]
For any valid \texttt{TemplateData}, the symmetric bilinear form on the $(3n+1)$-dimensional space is congruent to a diagonal matrix with exactly one positive and $3n$ negative entries. Formally:
\begin{lstlisting}
axiom basis_orthogonal_exists {n : Nat} (data : TemplateData n)
    (h_d : Rle data.d zero)
    (h_c_neg : True)
    (h_trace : True) :
    HasSignature (3 * n + 1) (templateMatrix data) 1 (3 * n)
\end{lstlisting}
\end{axiom}

\subsection{Axiom 2: Signature-to-Hodge Principle}
\begin{axiom}[Hodge Index Transfer]
If a space has signature $(1, q)$ and a vector $H$ has positive norm, then the restriction of the form to $H^\perp$ is negative semidefinite (signature $(0, q)$).
\begin{lstlisting}
axiom hodge_index_from_signature {dim : Nat} (M : Matrix dim dim) (H : Fin dim -> Real) (q : Nat)
    (h_sig : HasSignature dim M 1 q)
    (h_H_pos : Pos (eval M H H))
    (D : Fin dim -> Real)
    (h_orth : Req (eval M H D) zero) :
    Rle (eval M D D) zero
\end{lstlisting}
\end{axiom}

\subsection{Remark on Honesty}
These axioms are declared as trusted external facts. They are not proved in Lean. The existence of the arithmetic surface $\mathbb{S}$ is also unconstructed. The formal system enforces that these assumptions are auditable.

\section{Finite Computational Instantiation}
To validate the scaffold, we instantiate the template for the first seven primes \(P = \{2,3,5,7,11,13,17\}\), i.e., \(n=7\).

\subsection{Data Injection}
We generated high-precision rational intervals for $\log p$ (error $< 10^{-15}$) via an external Python script. The data is injected into \texttt{LefschetzData.lean} as an instance of \texttt{TemplateData}.

\subsection{Signature Verification}
The full-space Gram matrix (rank $\rho = 3n+1 = 22$) was computed. Applying the template theorem (conditional on the axioms) yields:
\[
\operatorname{Signature}(\text{Gram}) = (1, 21) = (1, \rho - 1).
\]
This matches the expected Hodge index signature for a surface. Consequently, via the Hodge index transfer axiom, we derive that the restriction to $H^\perp$ is negative definite for this finite truncation.

\begin{table}[h]
\centering
\begin{tabular}{|c|c|c|c|}
\hline
\textbf{Prime Set} & $\mathbf{n}$ & $\mathbf{\rho = 3n+1}$ & \textbf{Signature} \\
\hline
$\{2,3,5,7,11,13,17\}$ & 7 & 22 & $(1, 21)$ \\
\hline
\end{tabular}
\caption{Finite truncation signature verification.}
\label{tab:signature}
\end{table}

\section{Conclusion \& Future Directions}
We have presented a rigorous, machine-checked conditional proof: \textit{the existence of the F\textsubscript{1}-square with a Hodge index theorem implies the Riemann Hypothesis}. The formal scaffold includes a finite consistency check for the $T_1$ truncation.

\subsection{Open Problems}
\begin{enumerate}
\item \textbf{Proving the Linear Algebra Axioms:} The two axioms (orthogonal basis existence and signature-to-Hodge) are currently admitted. Proving them from first principles in constructive real analysis would strengthen the formalization.
\item \textbf{Generalizing the Induction:} Extending the proof to arbitrary $n$ (beyond the finite check) and deriving a finite-height RH statement.
\item \textbf{Constructing $\mathbb{S}$:} The ultimate goal remains the construction of $\operatorname{Spec} \mathbb{Z} \times_{\mathbb{F}_1} \operatorname{Spec} \mathbb{Z}$ with the required intersection theory.
\end{enumerate}

\subsection*{Disclaimer}
This document describes a research program. The Riemann Hypothesis remains open. The arithmetic surface $\mathbb{S}$ is unconstructed. No claim of an unconditional proof of RH is made.

\end{document}
